On one core of this laptop, every result reduces to the cache lines a query moves from the
level of the hierarchy the data lives in, and to what the CPU does with each line.

\paragraph{Layouts.} DSM wins whenever a query reads few attributes or produces few rows: up to
15$\times$ over NSM for one attribute of 16, and up to 9$\times$ over its own early strategy for
\code{SELECT *} at 0.01\%. NSM wins the reconstruction of many attributes over many rows,
because its reconstruction is a copy, and it is the least sensitive to strategy, width and
placement. PAX matches DSM only with large pages: with the 8--16\,KB pages of disk-based PAX its
minipages are too short for the prefetcher, and from 64\,KB on they alias in L1 unless their
stride is an odd number of lines. Where the data lives moves the strategy crossovers by an order
of magnitude, while the layout ranking stays the same.

\paragraph{Materialization.} Late reads, of every attribute it gathers, the fraction
$1-(1-s)^v$ of lines holding a selected value, and each gathered line costs 2.2--3.2 streamed
ones from DRAM. It wins below 2.4\% (4 attributes) to 7.4\% (16) from DRAM, and up to 20--33\%
for projections in L3, where gathered lines are cheap; in cache the CPU work per row decides,
and late SUM wins only with a hand-written SIMD selection vector. A threshold from the line size
and two fitted constants, chosen by where the table lives, picks between them per span and beats
both fixed strategies on clustered data.

\paragraph{Encodings.} Fewer bytes alone do not make a scan faster: decoding first was slower
than plain for every encoding, even rle with 32$\times$ fewer bytes. Compression wins when the
query works on the compressed form (SUM over runs 87$\times$, SUM from deltas 1.3$\times$, zone
maps up to 15$\times$ over late on plain, byte-aligned dictionary codes 2.2--5.4$\times$).
Otherwise it wins only if decoding costs less than the time memory takes per plain value, about
1 cycle from DRAM and 0.3 in cache: here that requires AVX2, and only pays from DRAM.

\paragraph{Tuning.} One hardware parameter decides one code choice: the line size and where the
data lives decide the materialization threshold, the cache sizes and the work per value decide
the vector size, the L1 set stride decides where columns and minipages start, write-allocate
decides how output is written, the prefetcher and the TLB decide the PAX and OS page sizes, and
the SIMD width and DRAM bandwidth decide the kernels and the code width. Together they make the
engine 1.46$\times$ faster over 18 query--layout pairs, up to 3.8$\times$, without changing the
queries.

\paragraph{Limitations.} All results come from one core of one laptop under WSL2, with 8-byte
integer attributes. The rules use parameters that sysfs reports, but the constants of the
threshold were fitted here, and the size of each effect depends on this machine; on another
machine they should be measured again with \code{make experiments}.
